\chapter{Implementation}

% =====================================================================
% NOTA DE USO
% Los títulos (\section / \subsection) están en inglés y son los definitivos.
% El texto en cursiva es SOLO una guía de contenido en español: bórralo al
% redactar.
%
% CRITERIO DE ALCANCE DE ESTE CAPÍTULO:
%   - SÍ: artefactos construidos, decisiones de realización, parámetros,
%     estructura del código, mecanismos que garantizan la validez del montaje.
%   - NO: qué son los frameworks (Background), por qué se eligió la política
%     de decisión y su formulación (Design), definición de los escenarios
%     evaluados y sus resultados (Evaluation).
% Regla de corte: si una frase seguiría siendo cierta aunque no se hubiera
% escrito una sola línea de código, no pertenece a este capítulo.
% =====================================================================

\emph{Párrafo introductorio: declarar explícitamente el alcance del capítulo.
Indicar que los frameworks de simulación se describen en el capítulo de
fundamentos y que la formulación de la política de decisión se desarrolla en el
capítulo de diseño, de modo que aquí se documenta únicamente su realización: qué
módulos se construyeron, qué código se extendió, qué se reutilizó sin modificar y
qué mecanismos garantizan que el montaje produce medidas válidas. Anticipar el
orden: primero la infraestructura, después el nodo y sus modelos, después la
aplicación y las capas de decisión, y por último la instrumentación y la
verificación del propio software.}

% =====================================================================
\section{Implementation Overview}
% =====================================================================

\subsection{Developed, Extended and Reused Components}
\emph{Tabla-resumen que abre el capítulo y lo hace navegable: cada componente
clasificado en desarrollado desde cero, extendido a partir de código existente o
reutilizado sin modificar, con la sección del capítulo donde se detalla. Es la
sección que delimita con precisión la contribución de ingeniería propia frente a
lo heredado de los frameworks.}

\subsection{Integration of the INET, Simu5G and Veins Code Bases}
\emph{Documentar los puntos de fricción concretos que hubo que resolver para que
las tres bases de código coexistieran en un mismo nodo y una misma compilación:
alineación de versiones, conflictos entre la pila IP de INET y la pila celular de
Simu5G, y uso de la movilidad externa dentro de un nodo celular. Solo problemas
reales y su solución, sin describir los frameworks.}

\subsection{Project Layout and Build Organisation}
\emph{Describir la organización del proyecto propio y su enlazado contra los tres
frameworks, la separación entre código, escenarios y resultados, y el criterio de
mantener acotadas y trazables las modificaciones realizadas sobre el código base
externo.}

% =====================================================================
\section{Mobility Scenario Construction}
% =====================================================================

\subsection{Road Network Generation from Cartographic Data}
\emph{Documentar el proceso seguido para obtener la red viaria del área de
estudio a partir de datos cartográficos abiertos: delimitación del área,
conversión al formato del simulador de tráfico, filtrado de tipos de vía y
depuración de la red resultante. Indicar las correcciones manuales necesarias.}

\subsection{Route, Stop and Traffic Demand Definition}
\emph{Describir la construcción de los ficheros de escenario: la ruta del
vehículo instrumentado, la declaración de sus paradas y la generación del tráfico
de fondo. Aquí se documenta cómo se generan los ficheros; la elección de qué
niveles de demanda se evalúan pertenece al capítulo de evaluación.}

\subsection{Building-Based Obstacle Generation}
\emph{Documentar la extracción de la geometría de edificios y su conversión al
formato de obstáculos consumido por el modelo de propagación, incluyendo el
procesamiento intermedio que fue necesario implementar.}

\subsection{Runtime Coupling and Exposed Mobility State}
\emph{Describir el acoplamiento en ejecución entre el simulador de tráfico y el
de red: sincronización temporal, instanciación dinámica de nodos y, sobre todo,
qué variables de movilidad quedan accesibles a los módulos de aplicación
(posición, velocidad, estado e identificador de parada). Esta última parte es la que
enlaza con la capa de predicción.}

% =====================================================================
\section{Hybrid Vehicular Node}
% =====================================================================

\subsection{Extension of the NR User Equipment with an IEEE 802.11 Interface}
\emph{Núcleo estructural del capítulo. Documentar el módulo compuesto construido
por extensión del equipo de usuario celular: submódulos añadidos, puertas de
radio declaradas, tipo de gestión inalámbrica empleado y conexión de la nueva
interfaz a la capa de red heredada. Incluir el diagrama del nodo y el fragmento
de definición del módulo.}

\subsection{Dual-Stack Addressing and Interface Binding}
\emph{Documentar la configuración que permite a un mismo nodo mantener
direcciones activas sobre ambas interfaces, y el mecanismo por el que se fuerza
que cada paquete salga por la interfaz que decide la aplicación y no por la que
seleccionaría el encaminamiento por defecto. Es un requisito sin el cual la
decisión de descarga sería inefectiva.}

\subsection{Infrastructure and Server-Side Modules}
\emph{Describir los módulos del lado fijo tal como quedaron configurados: puntos
de acceso asociados a las paradas, estación base celular, red de transporte y
servidor de recolección único. Indicar la ubicación de los elementos y los
parámetros de radio comunes.}

% =====================================================================
\section{Event-Driven Energy Consumption Model}
% =====================================================================

\subsection{Consumer Module Structure and Signal Subscription}
\emph{Documentar el módulo consumidor implementado: su punto de anclaje en la
jerarquía del nodo, el conjunto de señales de actividad a las que se suscribe y
el alcance de la suscripción. Enumerar las señales en una tabla.}

\subsection{Activity Classification and Current-State Transitions}
\emph{Describir la máquina de estados implementada: cómo cada señal se clasifica
en actividad de transmisión o recepción, cómo se resuelve la concurrencia entre
ambas, y el temporizador que devuelve el consumo al nivel base ante la ausencia
de actividad.}

\subsection{Energy-Domain Isolation in a Dual-Interface Node}
\emph{Sección crítica y de aportación propia. Documentar el problema detectado al
instalar el modelo energético celular en un nodo que también incorpora interfaz
inalámbrica —las señales de esta última activan el consumidor celular e inflan su
consumo— y la solución implementada para que cada dominio contabilice
exclusivamente su propio tráfico. Argumentar que sin esta corrección ninguna
comparación energética entre alternativas sería válida.}

\subsection{Battery Coupling and Recorded Energy Quantities}
\emph{Documentar la integración de carga en el tiempo, el acoplamiento con el
módulo de almacenamiento y las magnitudes finales registradas por cada dominio
energético, que son las que alimentan el capítulo de resultados.}

\subsection{Energy Parameter Set}
\emph{Presentar la tabla de parámetros efectivamente configurados en ambos
dominios (tensión de alimentación y corrientes de reposo, recepción y
transmisión), indicando su procedencia documental.}

% =====================================================================
\section{Hybrid Telemetry Application}
% =====================================================================

\subsection{Module Structure: Sender and Per-Technology Receivers}
\emph{Documentar la arquitectura implementada: emisor embarcado con acceso a
ambas interfaces y receptores diferenciados por tecnología en el servidor, con su
convergencia en un registro unificado de recepción. Incluir el diagrama de
módulos y sus interconexiones.}

\subsection{Message Definitions and Telemetry Metadata}
\emph{Documentar los formatos de mensaje definidos para cada categoría de datos,
sus tamaños lógicos, y los metadatos incorporados que permiten seguir cada
paquete extremo a extremo (identificadores, marcas de tiempo de generación,
interfaz de salida).}

\subsection{Finite Buffers and Overflow Handling}
\emph{Documentar la implementación de los buffers por categoría, la contabilidad
de ocupación en bytes, la política aplicada al saturarse y el registro de los
descartes producidos.}

\subsection{Interface Dispatch and Transmission Pacing}
\emph{Documentar cómo se materializa una decisión de encaminamiento sobre un
paquete concreto: selección de interfaz, extracción ordenada del buffer y
espaciado entre envíos, justificando este último por la saturación de colas
observada al enviar en ráfaga.}

\subsection{Acknowledgement, Retransmission and Duplicate Filtering}
\emph{Documentar la capa de fiabilidad implementada sobre datagramas:
confirmaciones agregadas, detección de pérdidas por número de secuencia,
retransmisión y filtrado de duplicados en recepción. Indicar los parámetros
configurables de esta capa.}

\subsection{Connectivity Transitions and In-Flight Data Recovery}
\emph{Documentar el tratamiento de la entrada y salida de cobertura: eventos de
asociación y desasociación, recuperación de los datos que quedaron en tránsito y
garantía de que ningún dato queda retenido en una interfaz ya no disponible.}

\subsection{Reliable-Transport Variant}
\emph{Documentar la variante de transporte fiable como adaptación mínima del
mismo diseño, precisando exactamente qué cambia y qué permanece idéntico, el
mecanismo de recuperación de sesión en los cambios de interfaz y la contabilidad
separada de la sobrecarga de tramado.}

% =====================================================================
\section{Prediction Layer Implementation}
% =====================================================================

\subsection{Offline Profile Extraction from Reference Runs}
\emph{Documentar el procedimiento implementado para construir los perfiles
históricos: ejecuciones de referencia del escenario, extracción de tiempos entre
paradas y de la duración efectiva de las ventanas de cobertura, y las
herramientas desarrolladas para ello.}

\subsection{Profile Data Format and Configured Percentiles}
\emph{Documentar el formato de los perfiles resultantes y los percentiles
efectivamente configurados para cada magnitud. La justificación del criterio de
percentiles corresponde al capítulo de diseño: aquí se indican los valores
adoptados y dónde se consumen.}

\subsection{Runtime Estimator Module}
\emph{Documentar el módulo que produce la estimación durante la simulación:
entradas de movilidad que consume, selección del siguiente punto de cobertura
válido, salidas expuestas al emisor y cadencia de actualización.}

\subsection{Live-Signal Arbitration and Fallback Behaviour}
\emph{Documentar dos mecanismos de robustez implementados: la prioridad de la
observación real de asociación sobre la predicción histórica cuando el vehículo
ya se encuentra en cobertura, y la degradación controlada ante perfiles o puntos
de cobertura ausentes, que emite advertencia sin detener la simulación.}

% =====================================================================
\section{Offloading Policy Implementation}
% =====================================================================

\subsection{Framework-Independent Decision Engine}
\emph{Documentar la decisión de realización más relevante de esta sección:
implementar el motor de decisión como componente sin dependencias del simulador,
lo que permite verificarlo de forma aislada y reutilizarlo sin modificación
dentro del emisor. Describir su interfaz de entrada y salida.}

\subsection{Category Registration and Per-Tick Input Assembly}
\emph{Documentar la generalización a un conjunto ordenado de categorías de datos
y cómo el emisor compone, en cada ciclo, el vector de entrada del motor a partir
del estado de los buffers y de la estimación vigente.}

\subsection{Discretisation of the Offloading Volume into Packets}
\emph{Documentar la conversión de la magnitud continua que produce la
formulación en un número entero de paquetes efectivamente enviables, y el
tratamiento del residuo. Es el punto donde la formulación teórica se vuelve
ejecutable.}

\subsection{Policy Execution Cycle and Drain Ordering in the Sender}
\emph{Documentar la ejecución periódica de la política dentro del emisor, su
independencia respecto de la cadencia de generación de datos, el orden de
prioridad entre categorías al vaciar y la ejecución efectiva del vaciado sobre
cada interfaz.}

\subsection{Deadline-Driven Rescue Layer}
\emph{Documentar la realización de la extensión por plazos: réplica de los plazos
en el emisor, evaluación de la condición de rescate en cada ciclo y las trazas
específicas que produce. La motivación y la formulación corresponden al capítulo
de diseño.}

\subsection{Precedence Between the Two Decision Layers}
\emph{Documentar cómo conviven ambas capas sobre el mismo emisor y qué regla de
precedencia se implementó cuando ambas concluyen simultáneamente, ya que de ello
depende el comportamiento observado en los resultados.}

\subsection{Boundary Cases and Safe Degradation}
\emph{Documentar los casos límite tratados explícitamente en el código: fin de
recorrido sin oportunidades futuras, arranque del sistema y ausencia de
predicción, argumentando que el sistema degrada de forma segura en todos ellos.}

% =====================================================================
\section{Configuration Management and Experimental Control}
% =====================================================================

\subsection{Single Base Configuration and Variant Derivation}
\emph{Documentar el mecanismo por el que todas las variantes se derivan por
herencia de una única configuración base, de modo que un cambio en el escenario
común alcanza simultáneamente a todas ellas y no puede producirse deriva entre
configuraciones. La definición de qué variantes se evalúan pertenece al capítulo
de evaluación.}

\subsection{Enforcing the Single-Difference Principle}
\emph{Documentar cómo la implementación garantiza que cada variante difiere de su
referencia en un único conmutador, y por qué esta propiedad —que es una decisión
de implementación, no de análisis— es la que hace comparables las mediciones
posteriores.}

\subsection{Exposed Parameters and Their Default Values}
\emph{Presentar en una tabla los parámetros que la implementación expone a la
experimentación, con su valor por defecto y el módulo que los consume. Sirve de
referencia cruzada para el capítulo de evaluación.}

% =====================================================================
\section{Instrumentation and Data Collection}
% =====================================================================

\subsection{Per-Packet and Per-Event Logging Infrastructure}
\emph{Documentar el sistema de trazas implementado y enumerar los registros
producidos: transmisión, recepción, desbordamiento, ocupación de buffer,
decisiones de la política, eventos de parada y de fiabilidad. Justificar la
necesidad de trazabilidad por paquete frente a los agregados nativos.}

\subsection{Native Statistics and Energy Scalars}
\emph{Documentar las magnitudes registradas mediante los mecanismos estadísticos
propios del simulador y de los modelos energéticos, cuyo papel es servir de
fuente independiente frente a los contadores de la aplicación.}

\subsection{Output Isolation Across Configuration Variants}
\emph{Documentar el esquema de nombrado de ficheros de salida que impide que una
variante sobrescriba los resultados de otra. Conviene declarar que responde a un
problema real detectado durante el desarrollo.}

\subsection{Post-Processing Tooling}
\emph{Documentar las herramientas desarrolladas para transformar las trazas
crudas en las métricas finales: agregación por ejecución, cálculo de las
magnitudes de entrega, retardo, frescura, desbordamiento y energía, y generación
de las tablas consumidas por el capítulo de resultados.}

% =====================================================================
\section{Verification of the Implementation}
% =====================================================================

\subsection{Stand-Alone Testing of the Decision Engine}
\emph{Documentar la verificación del motor de decisión fuera del simulador,
con casos construidos para cubrir umbrales y casos límite. Es lo que sostiene la
afirmación de que un comportamiento observado en los resultados procede del
escenario y no de un error de cálculo.}

\subsection{Offline Replay Validation of the Prediction Layer}
\emph{Documentar la validación de la capa de predicción por reproducción sobre
trazas reales, comparando la estimación producida con la ventana de cobertura
efectivamente observada, y el criterio de aceptación aplicado.}

\subsection{Consistency Checks Against Native Simulator Statistics}
\emph{Documentar la verificación cruzada entre los contadores propios de la
aplicación y las estadísticas nativas del simulador, y el umbral de discrepancia
por debajo del cual una ejecución se considera consistente.}

\subsection{Reproducibility: Seeds, Repetitions and Run Identification}
\emph{Documentar la gestión de semillas, el mecanismo de repeticiones y el
esquema de identificación de ejecuciones que permite reproducir exactamente
cualquier resultado presentado en la tesis.}

% =====================================================================
\section{Implementation Constraints and Deviations from the Design}
% =====================================================================
\emph{Sección de cierre breve y estrictamente acotada: recoger únicamente las
simplificaciones introducidas al implementar y los puntos en que la realización
se aparta del diseño formulado, indicando en cada caso por qué la desviación es
aceptable y qué efecto tiene sobre los resultados. Las limitaciones generales del
estudio corresponden al capítulo de conclusiones y no deben repetirse aquí.}
